\pdfoutput=1
\ifdefined\pdftrailerid\pdftrailerid{}\fi
\documentclass[11pt,a4paper]{article}
\usepackage[margin=25mm]{geometry}
\usepackage[T1]{fontenc}
\usepackage{lmodern,microtype}
\usepackage{amsmath,amssymb,amsthm,mathtools,booktabs}
\usepackage{xcolor,hyperref}
\hypersetup{colorlinks=true,linkcolor=blue!48!black,citecolor=blue!48!black,
  urlcolor=blue!52!black,pdftitle={Cycle logarithms and gauge discrimination in paired-cap sewing},pdfauthor={}}
\setlength{\emergencystretch}{2em}
\newtheorem{theorem}{Theorem}[section]
\newtheorem{proposition}[theorem]{Proposition}
\newtheorem{lemma}[theorem]{Lemma}
\newtheorem{corollary}[theorem]{Corollary}
\theoremstyle{remark}\newtheorem{remark}[theorem]{Remark}
\DeclareMathOperator{\coker}{coker}
\DeclareMathOperator{\Hom}{Hom}
\DeclareMathOperator{\sig}{sig}
\DeclareMathOperator{\diag}{diag}
\DeclareMathOperator{\rank}{rank}
\newcommand{\SU}{\mathrm{SU}}
\newcommand{\ZZ}{\mathbb Z}
\newcommand{\FF}{\mathbb F}
\newcommand{\RR}{\mathbb R}
\newcommand{\CC}{\mathbb C}
\newcommand{\dd}{\mathrm d}
\newcommand{\Ind}{\mathbf1}
\title{Cycle logarithms and gauge discrimination\\in paired-cap sewing}
\author{}
\date{}
\begin{document}
\maketitle
\begin{abstract}
We characterize two different kinds of information in closed
three-dimensional topological quantum amplitudes for a class of
paired-cap graph manifolds. On each fixed even cycle \(C_{2m}\), an exact
degree-two transfer converts the oriented \(\SU(2)\) invariant into a
positive color sum. Its common logarithmic scale gives
\(\binom{2m-2}{m-1}\pi^{-2m}r^{2m}\log r+O_m(r^{2m})\) in physical
normalization, along both parities. On a separate signed \(K_{4,4}\)
backbone, charge conjugation and vanishing sewing signature hide the seam
signs at all levels, even with block-interior decorations. The two
homologies differ only at the prime two, and no odd pointed modular
category or twisted finite odd nilpotent gauge theory distinguishes the
pair. Switching to the nonnilpotent odd gauge group
\(\FF_7\rtimes_{2}C_3\) does distinguish it, with the same distinction
for all three twists. An exact finite-field rank calculation determines
the two values. Group order \(21\) is minimal among odd-order finite
gauge groups for this pair.
\end{abstract}

\section{Information retained by a closed amplitude}

Closed partition functions are observables of a chosen theory, not
complete lists of the relations that define the underlying manifold.
The same paired-cap construction separates two questions. First, can a
closed \(\SU(2)\) scalar measure the size and shared scale of a genuine
backbone cycle? Second, when the graph and local blocks are fixed, can
different seam relations be invisible to that scalar and to whole
classes of additional quantum probes? We answer the first question on
even cycles and the second on a different, fixed complete bipartite
backbone. There is no inference from the alternating cycle estimate to
the asymptotics of the complete bipartite example.

The cycle proof uses the exact transfer and oriented normalization of
Papers~50 and~50B \cite{P50,P50B}. The new reduction is a single
translation coordinate in the logarithms of alternating colors; the
relative coordinates have an exponentially integrable closed-chain
weight. Its convolution trace fixes both the leading numerical
coefficient and the limiting positive folded-color law. This is a
closed uninserted amplitude, not a Wilson-angle distribution or a
semiclassical attribution to all Chern--Simons strata as studied on
other stocks in Paper~51 \cite{P51}.

The signed collision extends Paper~50B's odd-localized calculation to a
full integral relation matrix. On the selected pair, all odd-order
nilpotent representations forget the vertex groups because every
central quotient kills their images. The semidirect product
\(C_7\rtimes C_3\) instead retains a nontrivial transport character:
graph letters conjugate the \(C_7\) fibre variables in opposite
directions at the two ends of a seam. The resulting pair of twisted
biadjacency matrices, rather than ordinary abelian homology, shifts the
closed partition function. This is a change of gauge theory on the
same pair of manifolds, with no new surgery insertions. Known exact
quantum-invariant collisions (e.g. \cite{Funar}) motivate the question;
we do not claim that all-level collision or finite-group detection
itself is new.

\section{Common paired-cap sewing}\label{sec:sewing}

Let \(G=(V,E)\) be a finite simple connected graph with at least one
edge and degree \(d_v\) at \(v\). The unfilled vertex block is
\(X_v=\Sigma_{0,d_v+2}\times S^1\) with base-then-fibre orientation.
At every port \(z\) is the oriented base section and \(h\) is the
fibre, with \(z\cdot h=+1\). Fill two specified ports by the primitive
slopes \(p_vz+h\) and \(-p_vz+h\), with \(p_v\in\ZZ_{>0}\).
On the other ports, glue the two ends of edge \(e=vw\) by
\[
        z_{ve}\longmapsto \epsilon_e h_{we},\qquad
        h_v\longmapsto \epsilon_e z_{we},\qquad \epsilon_e\in\{+1,-1\}.
\]
Both maps reverse boundary orientation. Write \(M_{G,\epsilon,\mathbf p}\)
for this closed oriented manifold. Default cap orders are \(p_v=2\), and
then we suppress \(\mathbf p\). Write \(A_\epsilon\) for the signed
adjacency, \(n=|V|\), \(e=|E|\), and \(b_G=e-n+1\).

\subsection{Composition phase and the sign of a seam}

In the torus spaces use the section Lagrangian and the core basis
paired against its oriented dual. The basic sewing formula is the
Verlinde product-block contraction with the modular cap and Hopf seam
states \cite{BHMV,P50,P50B}. For a zero-weight lift preserving the
section Lagrangian, reversal of both torus coordinates is the
involution \(C=S^2\), charge conjugation. On an unpunctured torus this
can also be read by reversing the solid-torus core in the dual-paired
basis. Thus the change of sign of a seam inserts the identity in any
theory with self-dual simple labels. This statement is about the bare
contractions; it does not by itself determine their framing phase.

Here is the global phase check, including varying cap orders. Let
\(q,p\) be the section and fibre coordinates on the combined port
space, and let \(U\) take their per-vertex sums. On the block side the
homology kernel is
\(L_X=\{(U^{\mathsf T}a,p):Up=0\}\); on the paired cap and seam side
it is \(L_Y=\{(Tp,p)\}\). For the latter description, \(T\) has the cap
entries \(p_v,-p_v\) and the seam blocks
\(\left(\begin{smallmatrix}0&-\epsilon_e\\-\epsilon_e&0\end{smallmatrix}\right)\).
The two reciprocal cap entries cancel in the sum at each vertex.
Thus \(T\) is invertible, \(\sig T=0\), and \(UT^{-1}U^{\mathsf T}=-A_\epsilon\).
With the section reference \(L_0=\{(q,0)\}\), the cap/seam-to-block
Maslov form is the restriction of \(T\) to \(\ker U\)
\cite{GM,P50B}. This follows from writing the representative
\((Tp,p)\) as \((0,p)+(Tp,0)\) for \(Up=0\). Its signature is
\[
 \sig(T|_{\ker U})=\sig T-\sig(UT^{-1}U^{\mathsf T})=\sig A_\epsilon .
\]
For the degenerate case, the \(T\)-orthogonal complement of \(\ker U\)
is \(T^{-1}\operatorname{im}U^{\mathsf T}\). A joint radical of the two
restrictions and a dual isotropic space split off as a hyperbolic
summand of signature zero. This gives the same signature formula
without an invertibility assumption on \(A_\epsilon\).
On a bipartite graph \(\diag(+1,-1)\), with signs constant on the
parts, conjugates \(A_\epsilon\) to \(-A_\epsilon\), so the signature
and the entire composition anomaly vanish. In particular this
argument proves equality of \emph{oriented complex} scalars, not only
of absolute values. Interior ribbon graphs leave these boundary
homology kernels unchanged.

\subsection{Exact \texorpdfstring{\(\SU(2)\)}{SU(2)} weights for cap order two}\label{sec:weights}

Use shifted level \(r=k+2\ge3\) and
\[
 Z_r(S^2\times S^1)=1,\qquad Z_r(S^3)=S_{11},\qquad
 \tau_r(M)=Z_r(M)/S_{11}.
\]
For core labels \(\ell,j=1,\ldots,r-1\) the modular matrices and
anomaly unit are
\[
 S_{\ell j}=\sqrt{\frac2r}\sin\frac{\pi\ell j}{r},\qquad
 \Theta_{\ell\ell}=e^{\pi i(\ell^2-1)/(2r)},\qquad
 \kappa_r=e^{3\pi i/4-3\pi i/(2r)}.
\]
At vertex \(v\) the paired cap has weight
\(\big|(S\Theta^2 S)_{\ell1}\big|^2\) and the remaining local weight
is \(S_{1\ell}^{-d_v}\). A seam contributes \(S_{\ell j}\).
The two cap states may exchange roles according to the geometric
framing convention without changing their product. The closed
contraction includes a trace along a cycle, with no divide by its
state-space dimension.

For clarity, the parity in the cap weight can be checked directly.
Let \(G_r(j)=\sum_{u\bmod 2r}e^{\pi i(u^2+ju)/r}\). The cosine
difference formula for the sine product gives
\((S\Theta^2 S)_{\ell1}\), up to a unit phase, as
\([G_r(\ell-1)-G_r(\ell+1)]/(2r)\).
Translation gives \(G_r(j+2)=e^{-\pi i(j+1)/r}G_r(j)\).
In the squared modulus of the complete sum, Fourier orthogonality
retains only the two differences \(0,r\). Hence
\(\left|G_r(j)\right|^2=4r\) if \(r+j\) is even and is zero
otherwise. With the defect label \(a=r-\ell\), the cap weight
therefore equals \(4r^{-1}\cos^2(\pi a/(2r))\) on odd \(a\) and
vanishes for even \(a\).

Put
\[
 K(x)=\frac{\sin(\pi x)}{\pi x},\quad K(0)=1,\qquad
 F_d(x)=\cos^2(\pi x/2)\left(\frac{\pi x}{\sin\pi x}\right)^d .
\]
Converting to odd defects gives the exact zero-insertion formula
\begin{align}
 Z_r(M_{G,\epsilon})
 &=(-1)^{re}\kappa_r^{\sig A_\epsilon}
   2^{2n-e/2}\pi^{-e}r^{3e/2-n}I_r(G),\label{eq:finite}\\
 I_r(G)
 &= \sum_{\substack{1\le a_v<r\\ a_v\ {\rm odd}}}
       \prod_{v\in V}F_{d_v}(a_v/r)
       \prod_{vw\in E}K(a_v a_w/r).\notag
\end{align}
In fact the two odd defects at the ends of an edge give
\(\sin(\pi(r-a_v)(r-a_w)/r)=(-1)^r\sin(\pi a_va_w/r)\).
The ensuing \(a_v^{d_v}\pi^{d_v}/r^{d_v}\) precisely cancels the
corresponding factors introduced by \(F_{d_v}\).
The remaining powers follow from \(n\) paired caps \(4/r\),
\(e\) Hopf factors \(\sqrt{2/r}\), and \(2e\) occurrences of
\(S_{1a}^{-1}\) in the vertex powers. The complete oriented formula
agrees with \cite{P50,P50B}. The sign \((-1)^{re}\) appears from
the defect coordinates, not the choice of seam signing.

\section{Fixed even cycles have one logarithmic scale}\label{sec:cycles}

\begin{theorem}\label{thm:cycle}
Fix \(m\ge2\). On \(C_{2m}\) use the cap orders two and any seam
signing. With the normalization of \S\ref{sec:weights},
\[
 Z_r(M_{C_{2m},\epsilon})=
 \frac{\binom{2m-2}{m-1}}{\pi^{2m}}\, r^{2m}\log r
     +O_m(r^{2m}).
\]
There is no parity restriction on \(r\). In particular the leading
coefficient is positive and the leading logarithmic degree is one.
\end{theorem}

\begin{proof}
The adjacency signature is zero, and \(e=2m\) makes the parity
factor in \eqref{eq:finite} one. Label alternate defects \(a_i\)
and the intervening ones \(b_i\), with \(a_{m+1}=a_1\).
Put \(s_a=\min(a,r-a)\). The exact degree-two identity of
Paper~50 \cite{P50} is
\begin{equation}\label{eq:transfer}
 \sum_{\substack{1\le b<r\\ b\ {\rm odd}}}
 F_2(b/r)K(ab/r)K(cb/r)
 =\frac{r}{4ac}\min(s_a,s_c).
\end{equation}
An elementary proof also fixes the upper endpoint. The left
summand equals \(1/(4ac)\) times
\(\sin(2a\theta)\sin(2c\theta)/\sin^2\theta\), where
\(\theta=\pi b/(2r)\). On the odd grid \(a,c\) can be replaced
in the two sine numerators by \(s_a,s_c\), because
\(\sin(\pi b-2a\theta)=\sin(2a\theta)\).
Use \(\sin(2s\theta)/\sin\theta=2\sum_{j=1}^{s}
\cos((2j-1)\theta)\). On the full \(r\)-point midpoint grid
\(\theta=(b\pi)/(2r)\), \(b=1,3,\ldots,2r-1\), the cosine
modes of frequencies \(1,\ldots,r-1\) have norm squared \(r/2\)
and are mutually orthogonal by the finite geometric series.
Here \(2s_a-1,2s_c-1<r\). The full-grid sum of the sine product
ratios is \(2r\min(s_a,s_c)\); reflection about \(\pi/2\)
preserves the product, and the midpoint term, if present, is zero.
Take half and restore \(1/(4ac)\).

Each \(b_i\) can now be summed independently, leaving
\begin{align}
 Z_r(M_{C_{2m},\epsilon})
   &=\frac{2^m r^{2m}}{\pi^{2m}}T_m(r),\label{eq:Treln}\\
 T_m(r)
   &=\sum_{\substack{1\le a_i<r\\ a_i\ {\rm odd}}}
       \prod_{i=1}^m f_r(a_i)\min(s_{a_i},s_{a_{i+1}}),&
 f_r(a)&=\frac{\pi^2}{4r^2\sin^2(\pi a/(2r))}.\notag
\end{align}
Indeed \(I_r=(r/4)^m T_m\), because each \(a_i\) appears as
an endpoint of two transfers, so \(f_r(a)=a^{-2}F_2(a/r)\).
The reduced sum is positive termwise.

Let \(R=\lfloor r/2\rfloor\) and fold each label to \(s_a\), counting
the midpoint only once. Its weight is
\[
 w_r(s)=\sum_{\substack{1\le a<r,\ a\ {\rm odd}\\s_a=s}}f_r(a)
     =\frac{\Ind_{s\ {\rm odd}}}{s^2}+O(r^{-2}),\qquad
 0\le w_r(s)\le C s^{-2},\quad 1\le s\le R.
\]
The constants are uniform in parity: the lower preimage obeys the
first expansion by the bounded function
\(\csc^2 u-u^{-2}\) on \(0<u\le\pi/4\), while the upper preimage
has denominator \(4r^2\cos^2(\pi s/(2r))\) and contributes
\(O(r^{-2})\).

For \(s_{m+1}=s_1\), use the identity
\begin{equation}\label{eq:logkernel}
 \prod_i\frac{\min(s_i,s_{i+1})}{s_i^2}
  =\prod_i\frac1{s_i}
    \exp\!\left[-\frac12\sum_i|\log s_i-\log s_{i+1}|\right].
\end{equation}
For any fixed \(c>0\),
\(\sup_{a\ge1}\sum_{b\ge1}b^{-1}e^{-c|\log(b/a)|}<\infty\),
by splitting at \(a\). Holding one \(s_i\) fixed and deleting an
edge of the cycle, summation of the resulting path from its leaves
applies this bound repeatedly. In the telescoping replacement of the
\(m\) factors \(w_r(s_i)\), the erroneous site's cost relative to
\(s_i^{-2}\) is bounded by \(C r^{-2}s_i^2\), while the other
weights are bounded by constants times \(s_j^{-2}\). Therefore
each error term sums to \(O_m(r^{-2}\sum_{s\le R}s)=O_m(1)\).
This is an absolute error bound for the entire folded measure.

Compare the remaining sum on odd anchors with the integral of
the right side of \eqref{eq:logkernel} over cells
\(\prod_i[a_i,a_i+2]\). Over such a cell the logarithm of that
integrand changes by at most \(C_m\sum_i a_i^{-1}\).
Its ratio to the corner value is uniformly bounded. Use the same
path estimate for each inserted \(a_i^{-1}\) and
\(\sum_{a_i\ge1}a_i^{-2}<\infty\) to bound the sum-integral
error by \(O_m(1)\). With \(t_i=\log s_i\), this gives
\[
 T_m(r)=2^{-m}\int_{[0,L]^m}
          e^{-\frac12\sum_i|t_i-t_{i+1}|}\,\dd^m t+O_m(1),
 \qquad L=\log(R+O(1))=\log r+O(1).
\]
The factor \(2^{-m}\) is the exact leading density of the odd
lattice, not a fixed-parity assumption about the level.

Fix relative coordinates \(v_i=t_i-t_1\), \(v_1=v_{m+1}=0\).
Write \(D(v)=\max_i v_i-\min_i v_i\). For given \(v\), \(t_1\)
ranges on an interval of length \((L-D(v))_+\). Since the closed
variation \(\sum_i|v_i-v_{i+1}|\) is at least \(2D(v)\), the relative
weight has finite mass and finite first \(D\)-moment. Thus the
integral is \(L C_m+O_m(1)\), where
\[
 C_m=\int_{\RR^{m-1}} e^{-\frac12\sum_i|v_i-v_{i+1}|}\,\dd^{m-1}v .
\]
Equivalently \(C_m=k^{*m}(0)\), \(k(x)=e^{-|x|/2}\). In the
Fourier convention \(\widehat k(\xi)=\int e^{-i\xi x}k(x)\,\dd x\)
the transform is \((\xi^2+1/4)^{-1}\), so the beta integral gives
\[
 C_m=\frac1{2\pi}\int_{\RR}(\xi^2+1/4)^{-m}\,\dd\xi
 =\frac{2^{2m-2}\Gamma(m-\frac12)}{\sqrt\pi\,\Gamma(m)}
 =\binom{2m-2}{m-1}.
\]
Substituting \(T_m=2^{-m}C_m\log r+O_m(1)\) into
\eqref{eq:Treln} proves the assertion.
\end{proof}

\begin{corollary}[Common exponent and relative log colors]\label{cor:law}
Normalize the folded positive weight in \(T_m(r)\) to probability.
Then \((\log s_1/\log r,\,(\log(s_i/s_1))_{i=2}^m)\) converges
weakly to the product of the uniform law on \([0,1]\) and the density
\[
 C_m^{-1}\exp\!\left[-\tfrac12\sum_i|v_i-v_{i+1}|\right]\dd^{m-1}v.
\]
Thus all the retained alternating labels share the same limiting
exponent, whereas their relative logs retain a nontrivial law.
\end{corollary}
\begin{proof}
The weight replacement is \(O_m(1)\) in total variation. Away from
large relative differences, discard \(t_1\le M\) at cost \(O_m(M+1)\)
in the unnormalized weights and compare the remaining small cells
with the integral as before. With \(\max|v_i|\) bounded, their
diameters in the log variables tend uniformly to zero as
\(M\to\infty\), so this comparison also applies to bounded continuous
tests on fixed compact relative sets. The removed part vanishes
after division by \(\log r\) for fixed \(M\).
In the integral, set \(u=t_1/\log r\); its allowed interval tends to
\([0,1]\) uniformly for bounded \(v\).
Both discrete and continuous relative tails are uniformly
exponentially small after normalization: the factor
\(e^{-\sum|v_i-v_{i+1}|/4}\) bounds the cost of the other half of
the exponent whenever \(\max|v_i|\) is large, and the remaining
half is summed by the same fixed-anchor path bound. Let the
relative cutoff then tend to infinity.
\end{proof}

This law is a probability interpretation of the exact reduced
uninserted scalar. It does not assume that the folded labels in the
proof are jointly measurable Wilson angles. Even-cycle length is
visible here through the growth exponent \(2m\); no assertion compares
that exponent with every other type of backbone.

\section{Integral seam relations and exact collisions}\label{sec:collision}

\subsection{What the closed torus scalar discards}

\begin{proposition}\label{prop:blind}
For fixed positive orders \(\mathbf p\) on a bipartite backbone,
every seam signing gives the same oriented \(\SU(2)\) RT scalar
at every shifted level \(r\ge3\). The equality remains true after
placing identical colored framed ribbon graphs wholly in the
interiors of corresponding capped blocks. In both statements
one may replace \(\SU(2)\) by a modular category with all simple
labels self-dual, in the extended-TQFT convention of
\S\ref{sec:sewing}.
\end{proposition}
\begin{proof}
Decorated blocks, like empty blocks, give tensors in unpunctured
torus spaces when no ribbon meets a seam. On this space \(-I\)
acts by charge conjugation and is the identity under the stated
hypothesis. The cap reciprocals cancel in the oriented Maslov
calculation of \S\ref{sec:sewing} for all paired orders \(\mathbf p\),
and the signed adjacency has zero signature by bipartiteness.
Thus neither the bare contraction nor the anomaly factor changes.
\end{proof}

For a ribbon crossing a sewing torus the relevant torus is punctured;
the proof gives no equality for that experiment. In particular it
does not turn all colored observables into functions of the
unpunctured contraction.

\begin{theorem}[Integral relation matrix]\label{thm:integral}
Let \(P=\diag(p_v)\). With the construction of \S\ref{sec:sewing}
there is a noncanonical isomorphism
\[
 H_1(M_{G,\epsilon,\mathbf p};\ZZ)
  \simeq \ZZ^{b_G}\oplus\coker_{\ZZ}(P A_\epsilon P).
\]
For paired order two the matrix is \(4A_\epsilon\).
\end{theorem}
\begin{proof}
Before the seams let \(a_v,b_v\) denote the sections of the
two capped ports. The integral relations at the block are
\[
 p_va_v+h_v=0,\qquad -p_vb_v+h_v=0,\qquad
 a_v+b_v+\sum_{e\ni v}z_{ve}=0.
\]
Eliminate \(h_v\) from the first and \(b_v\) from the last, using
unit coefficients. The second is then \(p_v\sum_e z_{ve}=0\).
The seam relations give \(z_{ve}=-\epsilon_e p_w a_w\).
Thus the remaining vertex/port cokernel is exactly that of
\(P A_\epsilon P\), without dividing by \(p_v\).
The next term in the graph-of-spaces Mayer--Vietoris sequence
is the kernel of the ordinary vertex-edge incidence map on
\(H_0\). This is \(\ZZ^{b_G}\), since seams induce the same
\(H_0\) maps for either sign. The extension by this free quotient
splits.
\end{proof}

\subsection{All odd pointed theories see only odd-order homomorphisms}

\begin{theorem}\label{thm:pointed}
For paired order-two caps on a bipartite graph, let \(B\) be the
signed biadjacency, with \(n_A\) rows and \(n_B\) columns.
For any finite nondegenerate odd-order metric group
\((\mathcal A,Q)\) of order \(m_0\), its vacuum-divided pointed RT
amplitude is
\[
 \mathcal P_{\mathcal A,Q}(M_{G,\epsilon})
 =m_0^{(e+1)/2-n_B}\,|\ker(B:\mathcal A^{n_B}\to\mathcal A^{n_A})|
 =\sqrt{|\Hom(H_1(M_{G,\epsilon};\ZZ),\mathcal A)|}.
\]
It is positive real and does not depend on \(Q\).
\end{theorem}
\begin{proof}
Write \(\beta(x,y)=Q(x+y)/(Q(x)Q(y))\). The torus matrices are
\(S_{xy}=m_0^{-1/2}\beta(x,y)\), up to conjugation, and
\(\Theta_{xx}=Q(x)\). A cap coefficient is
\(c_+(a)=m_0^{-1}\sum_x\beta(a,x)Q(x)^2\).
The opposite cap gives its conjugate, by \(x\mapsto-x\) and
\(Q(-x)=Q(x)\). Expanding the squared modulus with \(x=y+c\)
leaves
\[
 |c_+(a)|^2=m_0^{-2}\sum_c \beta(a,c)Q(c)^2
                                \sum_y\beta(y,c)^2=m_0^{-1}.
\]
Indeed the inner sum vanishes unless \(2c=0\) by nondegeneracy,
and odd order forces \(c=0\). Every vertex then contributes
\(m_0^{d_v/2-1}\). A seam contributes
\(m_0^{-1/2}\beta(x_i,y_j)^{\epsilon_{ij}}\), or the simultaneous
conjugate convention. The Maslov signature is zero. Dividing by
\(Z(S^3)=m_0^{-1/2}\) and summing first over the row colors gives
\[
  \mathcal P=m_0^{(e+1)/2-n_A-n_B}
    \sum_{x,y}\prod_{i,j}\beta(x_i,y_j)^{B_{ij}}
  =m_0^{(e+1)/2-n_B}|\ker B|.
\]
All vanishing and nonvanishing in this sum are exact character
orthogonality. Finally \(4\) is invertible on \(\mathcal A\);
Theorem~\ref{thm:integral} gives
\(|\Hom(H_1,\mathcal A)|=m_0^{b_G}|\ker B||\ker B^{\mathsf T}|\).
Perfect pairings give \(|\ker B^{\mathsf T}|=
m_0^{n_A-n_B}|\ker B|\), also if \(B\) is rectangular.
The square root is the positive one by the preceding calculation.
\end{proof}

\subsection{A concrete pair, and a first abelian discriminator}

For \(G=K_{N,N}\) use the signed biadjacency
\[
 B(D)=\begin{pmatrix}1&\mathbf1^{\mathsf T}\\
            \mathbf1&\mathbf1\mathbf1^{\mathsf T}-2D\end{pmatrix},
\]
where \(D\) is a nonsingular binary integer matrix of size \(N-1\).
Subtract the first row and column to put
\(B(D)\sim_{\ZZ}\diag(1,-2D)\). A permutation puts
\(A_\epsilon\) into \(B(D)\oplus B(D)^{\mathsf T}\), in the sense
of integer presentation equivalence. Thus, if \(d_i\) are the
positive Smith factors of \(D\), the group is
\[
 H_1(M_{B(D)};\ZZ)\simeq
 \ZZ^{(N-1)^2}\oplus(\ZZ/4)^2\oplus
 \bigoplus_{i=1}^{N-1}(\ZZ/(8d_i))^2 .
\]

Fix \(N=4\) and compare \(D_0=I_3\) and \(D_1=J_3-I_3\), with
\(J_3\) the all-ones matrix. Both have a \(2\times2\) minor of
modulus one and their determinant moduli are \(1\) and \(2\),
so their Smith lists are \((1,1,1)\) and \((1,1,2)\). The two
closed manifolds \(M_0,M_1\) have
\begin{align}
 H_1(M_0;\ZZ)&\simeq\ZZ^9\oplus(\ZZ/4)^2\oplus(\ZZ/8)^6,\notag\\
 H_1(M_1;\ZZ)&\simeq\ZZ^9\oplus(\ZZ/4)^2\oplus(\ZZ/8)^4
                                      \oplus(\ZZ/16)^2.\label{eq:groups}
\end{align}
They are not homeomorphic. Proposition~\ref{prop:blind} gives
the all-level oriented and block-interior decorated collisions,
and Theorem~\ref{thm:pointed} gives the common scalar
\(|\mathcal A|^{9/2}\) for every odd pointed modular theory.

As an elementary detector outside that class, for the untwisted
cyclic theory of order \(2^a\), use its \(S^3\)-divided partition
function \(|\Hom(H_1,\ZZ/2^a)|\). The groups in
\eqref{eq:groups} agree for \(a=1,2,3\), and give \(2^{58},2^{60}\)
at \(a=4\). This even-order abelian separator and the minimal
odd-order separator below are independent statements.

\section{The smallest odd-order finite-gauge separator}\label{sec:sep}

Let \(K\) be finite and \(\omega\in H^3(BK;U(1))\).
For a connected closed oriented three-manifold \(M\), our finite-gauge
normalization is
\begin{equation}\label{eq:dw}
 \mathcal D_{K,\omega}(M):=\frac{Z_{K,\omega}(M)}{Z_{K,\omega}(S^3)}
 =\sum_{\rho\in\Hom(\pi_1 M,K)}
      \big\langle (B\rho)^*\omega,[M]\big\rangle .
\end{equation}
Here \(B\rho\) denotes the induced classifying map from \(M\).
The physical sum uses weight \(1/|K|\) before vacuum division
\cite{DW,FQ}; equivalently sum once per conjugacy orbit with
reciprocal centralizer weight. No division by \(K\) is made in the
right side of \eqref{eq:dw}. Every weight in the results below is one;
in particular, the two computed physical amplitudes are positive.

\subsection{Why no odd nilpotent group can detect the pair}

\begin{theorem}\label{thm:nilp}
Suppose \(p_v=2\), \(\det A_\epsilon\ne0\) and \(K\) is finite
nilpotent of odd order, coprime to \(\det A_\epsilon\).
Let \(\pi:\pi_1 M_{G,\epsilon}\twoheadrightarrow F_{b_G}\) be
the natural graph quotient. Every \(K\)-representation factors
uniquely through \(\pi\). Hence for all \(\omega\),
\[
 Z_{K,\omega}(M_{G,\epsilon})=|K|^{b_G-1},\qquad
 \mathcal D_{K,\omega}(M_{G,\epsilon})=|K|^{b_G}.
\]
The pro-\(\ell\) completion is free pro-\(\ell\) of rank \(b_G\)
for every odd prime \(\ell\nmid\det A_\epsilon\).
\end{theorem}
\begin{proof}
Choose a maximal spanning tree and based port loops with the
two cap loops consecutive in the boundary product. The graph
presentation has a fibre \(h_v\) and cap sections \(a_v,b_v\)
at each vertex, with \(a_v^2=h_v^{-1}\), \(b_v^2=h_v\).
They are the only generators in addition to the ordinary seam
sections and the \(b_G\) stable letters. In a finite group of odd
exponent \(d\) squaring is bijective, with inverse the power
\(u\) for \(2u\equiv1\bmod d\). Consequently the images of \(a_v,b_v\)
are \(h_v^{-u},h_v^u\); they multiply to one and all the required
commutations with \(h_v\) hold in every such image.
The surviving vertex equation is an ordered product of its seam
section images equal to one. Each such section is a conjugate of
\(h_w^{\epsilon_{vw}}\), by the identity on tree seams and by a
stable letter or its inverse otherwise.

Let \(\gamma_1K=K\), \(\gamma_{j+1}K=[K,\gamma_jK]\).
If all fibre images lie in \(\gamma_jK\), conjugations have no
effect in \(\gamma_jK/\gamma_{j+1}K\). In this finite abelian
quotient the vertex equations are
\(A_\epsilon\overline{\mathbf h}=0\).
The adjugate annihilates \(\overline{\mathbf h}\) by
\(\det A_\epsilon\), which is invertible on the quotient by
coprimality. Thus the fibres lie in \(\gamma_{j+1}K\). Beginning
at \(j=1\), nilpotence forces all fibres, then all vertex groups,
to have trivial image. The remaining representation is arbitrary
on the free stable letters.

The projecting map is realized topologically by collapsing the
capped blocks and using intervals on the seam collars. Thus the
classifying maps in this calculation factor through a graph,
where degree-three cohomology vanishes. Formula \eqref{eq:dw}
now gives the stated numbers. The same argument applies to all
finite \(\ell\)-groups for the indicated primes, and identifies
the corresponding pro-\(\ell\) completions by their universal
property.
\end{proof}

For the pair \(M_0,M_1\), the determinant moduli of the signed
adjacencies are \(2^6\) and \(2^8\). Thus all finite odd
nilpotent theories, including all their twists, have the
common vacuum-divided value \(|K|^9\). This conclusion concerns
a whole finite-group class, not just abelianization.

\subsection{A conjugation channel at order twenty-one}

\begin{theorem}\label{thm:sep}
Let \(K=\FF_7\rtimes C_3\), with \(C_3\) acting on the additive
\(\FF_7\) by multiplication by \(2\). For every
\(\omega\in H^3(BK;U(1))\),
\[
 \mathcal D_{K,\omega}(M_0)=2{,}405{,}357{,}452{,}449,\qquad
 \mathcal D_{K,\omega}(M_1)=2{,}727{,}379{,}236{,}309 .
\]
There are three twists and each gives these same two values.
No odd-order finite gauge group of smaller order distinguishes
this pair, with any twist.
\end{theorem}
\begin{proof}
Write the multiplication as
\((x,e)(y,f)=(x+2^e y,e+f)\), with \(e,f\in\ZZ/3\).
Under the projection to \(C_3\) every representation kills all
vertex groups by Theorem~\ref{thm:nilp}; in particular all fibre
images lie in \(\FF_7\). The argument only uses invertibility of
the signed adjacency modulo three for this abelian quotient.
The cap images are then uniquely determined by their fibres,
and both the cap and seam sections lie in the abelian normal
subgroup. A stable letter \((u,e)\) conjugates a fibre by \(2^e\),
independently of \(u\).

Index each part of \(K_{4,4}\) by \(0,1,2,3\). Use the spanning
tree consisting of row zero and column zero. Tree exponents
are zero; the other nine exponents \(e_{ij}\) are independently
in \(\ZZ/3\). With \(b_{ij}=B(D)_{ij}\in\{+1,-1\}\), the fibre
relations, written additively in \(\FF_7\), are
\begin{equation}\label{eq:rank}
  A^+(\mathbf e)\,h_B=0,\qquad
  A^-(\mathbf e)^{\mathsf T} h_A=0,\quad
  A^+_{ij}=b_{ij}2^{e_{ij}},\quad
  A^-_{ij}=b_{ij}2^{-e_{ij}} .
\end{equation}
The reverse power comes from the inverse stable letter at the
other end of the same seam. The signs are the exponents
\(\epsilon_{ij}\) on the adjacent fibre, not new choices of
stable letters. After specifying both fibre vectors, all vertex
relations and every seam relation hold: the defined cap roots
commute with their fibres and the vertex products in the normal
subgroup are precisely the two systems in \eqref{eq:rank}.
There are no additional stable-letter restrictions, by the
graph-of-groups presentation. Hence the number of representations
at fixed exponents is
\(7^9\,7^\nu\), where
\(\nu=8-\rank_{\FF_7} A^+-\rank_{\FF_7}A^-\).
The factor \(7^9\) is from the independent normal components
of the nine stable letters. These are raw based homomorphism
counts, not counts modulo conjugation.

Here is the complete finite enumeration. Both first rows and
columns equal one, so subtracting their pivot reduces each
rank to one plus the rank of the \(3\times3\) matrix with entries
\(b_{ij}2^{\pm e_{ij}}-1\) for \(i,j\ge1\). For a \(3\times3\)
matrix the rank is \(3\) if its determinant is nonzero, else
\(2\) if a \(2\times2\) minor survives, else \(1\) if an entry
survives, else \(0\). Enumerate all nine exponents by the Cartesian
product of three choices, in row-major order. The following
histogram lists all nonempty bins:
\begin{center}
\begin{tabular}{lrrrr}
\toprule
seam matrix & \(\nu=0\) & \(\nu=1\) & \(\nu=2\) &
          \(\sum_{\mathbf e}7^\nu\)\\
\midrule
\(B(I_3)\)       & 14415 & 5070 & 198 & 59607\\
\(B(J_3-I_3)\)   & 13687 & 5712 & 284 & 67587\\
\bottomrule
\end{tabular}
\end{center}
For explicit reproducibility the entire arithmetic in each loop
is
\[
 (u^+,u^-)=(1,1),(2,4),(4,2)\quad
 \longmapsto\quad
 (b_{ij}u^+-1,b_{ij}u^--1)\pmod 7
\]
entry by entry. Start each histogram at zero; for each array of
these nine choices increment index
\(6-\rank(M^+)-\rank(M^-)\) once.
The accompanying script additionally verifies every rank by
ordinary modular elimination on the original \(4\times4\)
matrices. Each row sums to \(3^9=19683\), and the last column is
the respective dot product with \((1,7,49)\). Multiplication by
\(7^9\) gives the displayed exact vacuum-divided values in the
untwisted case. This finite enumeration is part of the calculation,
not an asymptotic experiment.

To address twists, consider \(1\to C_7\to K\to C_3\to1\).
For \(q>0\), \(H^q(C_7;U(1))\) is \(7\)-primary, so its
positive-degree \(C_3\) cohomology is zero by averaging.
The \(q=3\) group is \(\ZZ/7\) and the automorphism acts as
\(2^2=4\) (or its inverse for the opposite pullback convention).
Indeed \(H^3(C_7;U(1))\cong H^4(C_7;\ZZ)\) is generated by
the square of the degree-two cyclic class. There are no invariants.
Thus in total degree three the Lyndon--Hochschild--Serre terms
leave only the inflation from \(H^3(C_3;U(1))\cong\ZZ/3\).
The section \(C_3\hookrightarrow K\) makes inflation injective,
so it supplies all three twists. Each projected \(C_3\)
character factors through the backbone; the pullback of its
degree-three class therefore vanishes on the manifold.
All weights in \eqref{eq:dw} are exactly one for every twist.

Finally the positive odd integers below \(21\) are \(1\),
the primes below \(21\), \(9\) and \(15\). Groups of these
orders are nilpotent: order \(9\) is a prime power; at order
\(15\) both Sylow subgroups are normal by their counting
congruences, so the group is cyclic. Theorem~\ref{thm:nilp}
therefore applies to every smaller odd gauge group for this
pair, proving minimality.
\end{proof}

The mechanism in \eqref{eq:rank} is a fixed normal-module
character that survives conjugation by graph letters.
The lower central filtration of an odd nilpotent target instead
kills it. The result is not a classification of all nonnilpotent
groups or of detectors across other seam signings.


\section{Conclusions and scope}

In this construction the two effects have different mathematical
origins. The cycle sees an integrable distribution of relative
logarithmic colors and one noncompact common scale. The collision is
an exact insensitivity of the unpunctured torus contraction together
with a zero anomaly, even though the integer seam relations changed.
Odd pointed and nilpotent gauge probes cannot recover those relations
on the selected pair; coprime nonnilpotent conjugation can. Neither
example is a reconstruction of a general graph from closed invariants.
The fixed-cycle error is not uniform in cycle length, and the
decoration claim does not extend to graphs of ribbons crossing a
sewing torus. No all-orders semiclassical or marked Wilson theorem is
needed.

\begin{thebibliography}{99}
\bibitem{BHMV} C. Blanchet, N. Habegger, G. Masbaum and P. Vogel,
\emph{Topological quantum field theories derived from the Kauffman bracket},
Topology \textbf{34} (1995), 883--927.
\href{https://webhomes.maths.ed.ac.uk/~v1ranick/papers/bhmv.pdf}{Primary text.}
\bibitem{GM} P. Gilmer and G. Masbaum,
\emph{Maslov index, Lagrangians, mapping class groups and TQFT},
Forum Math. \textbf{25} (2013), 1067--1106.
\href{https://arxiv.org/abs/0912.4706}{arXiv:0912.4706.}
\bibitem{DW} R. Dijkgraaf and E. Witten,
\emph{Topological gauge theories and group cohomology},
Commun. Math. Phys. \textbf{129} (1990), 393--429.
\href{https://www.ias.edu/sites/default/files/sns/\%5B126\%5DCommMathPhys129-1990.pdf}
{Primary text.}
\bibitem{FQ} D. S. Freed and F. Quinn,
\emph{Chern--Simons theory with finite gauge group},
Commun. Math. Phys. \textbf{156} (1993), 435--472.
\href{https://arxiv.org/abs/hep-th/9111004}{arXiv:hep-th/9111004.}
\bibitem{Funar} L. Funar, with an appendix by L. Funar and A. Rapinchuk,
\emph{Torus bundles not distinguished by TQFT invariants},
Geom. Topol. \textbf{17} (2013), 2289--2344.
\href{https://arxiv.org/abs/1101.0509}{arXiv:1101.0509.}
\bibitem{P50} AI Automated Quantum Gravity Research,
\emph{What graph structure survives WRT sewing? Hidden connector
motifs and the stability boundary of dense cores}, Paper~50 (2026).
\href{https://github.com/g3173028019-beep/ai-automated-quantum-gravity-research/tree/main/papers/50-wrt-graph-detection-and-stability}{Source.}
\bibitem{P50B} AI Automated Quantum Gravity Research,
\emph{All-level WRT collisions on fixed graph-manifold backbones},
Paper~50B (2026).
\href{https://github.com/g3173028019-beep/ai-automated-quantum-gravity-research/tree/main/papers/50b-wrt-topological-blindness}{Source.}
\bibitem{P51} AI Automated Quantum Gravity Research,
\emph{Singular semiclassical sewing for paired-cap
Witten--Reshetikhin--Turaev invariants}, Paper~51 (2026).
\href{https://github.com/g3173028019-beep/ai-automated-quantum-gravity-research/tree/main/papers/51-singular-wrt-semiclassical-sewing}{Source.}
\end{thebibliography}
\end{document}
